\subsection{从公式回到像素}

先取一个方便手算的例子，只看红色通道。设投影边长 $w=10$，纹理坐标半径 $R=0.2$，墙面收到的光 $q=2$，反射率 $\rho_y=0.5$，于是光源图上存的是 $D_y=1$。假设整个圆盘上朝向和距离合起来给出 $G=0.25$。代入上页公式，白色表面的结果为 $(10\times0.2)^2\times 0.25/\pi\approx0.31831$。地面反射率为 $0.8$ 时，朝相机的结果约为 $0.25465$。若墙、地面和灯的覆盖范围一起放大两倍，圆盘面积变成四倍，$G$ 则降为四分之一，两者相抵。随文的 \path{verify.py} 用这组数复核系数、投影面积和尺度。

显卡测试使用另一组容易判断的输入。\path{tests/Renderer3DTests.cpp} 让 RSM 纹理处处返回 $(1,0.25,0)$，两面正对、相隔 1，投影边长为 4，纹理坐标半径为 0.2。此时红色通道应为 $0.8^2/\pi\approx0.20372$，OpenGL 实际读回与公式相差不到 0.001。只把距离改为 2，亮度降为原来的四分之一。反过来改变投影宽度和纹理坐标半径、让两者乘积保持相同，结果也相同。

最后回看 Cornell 的地面 A。最早的参考图把方向光照向绿墙，光源图里红墙的反射通量为零，调高间接光强度也不会让这堵墙染红地板。把灯转向红墙，再固定同一相机、灯光和强度 5 做算法前后对照，A 点的间接增量从 $(0.07715,0.03076,0.03076)$ 变成 $(0.48218,0.13550,0.13550)$。本节开头的展示图把强度调回 1，A 点的红色增量为 0.09644。先查灯有没有照到墙，再在固定强度下比较同一块地面，就能分清灯光方向和代码改动各自带来的变化。原始读数见 \path{output/scenes/rsm/measurements.csv}。

\subsection{仍会出现的误差}

收集阶段默认墙上的样本能照到地面。若在两点之间放一块挡板，程序仍可能把红光算过去。圆盘之外的反射面则完全收不到。靠近墙脚时，把最小距离从旧版固定的 \code{DIST\_BIAS=0.05} 与几何因子上限 4，改为随纹素宽度取值，能抑制很近的亮点。光源分辨率改变后，这一近似仍会跟着变。\source{https://doi.org/10.1145/1053427.1053460}{RSM 原论文第 3.3 节讨论了遮挡近似}

目前只有第一盏方向光参与这一次间接漫反射。墙面用几何表面的插值法线，材质的法线贴图没有进入这一步。这次计算只考虑漫反射，所以纯金属既不把漫反射光送到地面，也不接收这样的光，金属的镜面反射、多次反弹和材质随角度变化的完整 Fresnel 反射都在这次估算之外。两个独立绘制测试分别检查了受光面正照与斜照、背光面与金属，以及给金属接收面加减红色输入后的结果。

相机移动时，程序把当前点投回上一帧来平滑噪点。已有的 61 项场景检查包含 16、32、128 个圆盘样本以及切换投影、光源、场景和 RSM 设置后的历史重置。普通相机移动保留历史，上一帧裁剪范围以外的点会被拒绝。同一像素上一帧是否被其他物体挡住、物体自身移动到了哪里，仍缺少对应的深度和运动数据。测试代码见 \path{tests/DepthPrepassTests.cpp} 与 \path{tests/SceneRenderTests.cpp}，运行记录见 \path{output/validation/rsm-regression.txt}。
